% q0052.tex — Timing Is (Almost) Everything
\question{
  id        = {0052},
  title     = {Timing Is (Almost) Everything},
  source    = {OBECON},
  year      = {2026},
  round     = {Seletiva IEO -- Phase C},
  language  = {English},
  topic     = {Finance},
  subtopic  = {Valuation / Mid-Year Convention},
  type      = {objective},
  statement = {%
    A mining company generates continuous cash flows throughout each year.
    Analyst~A argues for the year-end convention (discounting all cash flows
    as if received at year-end); Analyst~B argues for the mid-year convention
    (discounting as if received at mid-year). Which approach is more accurate
    for valuation?
  },
  choices   = {%
    \item Analyst~A; year-end convention is standard practice regardless of
          actual cash flow timing.
    \item Analyst~B; mid-year convention better reflects continuous cash flow
          timing.
    \item Analyst~B; even though cash should be modelled as arriving at the
          beginning of each period.
    \item Both approaches are equally accurate as they discount the same cash
          flows.
  },
  answer    = {B},
  solution  = {%
    For evenly distributed continuous cash flows, the average timing is
    mid-year (month~6), not year-end (month~12). Using the year-end convention
    systematically undervalues the asset by assuming cash arrives on average
    6~months later than it actually does. The mid-year convention corrects
    this bias and produces more accurate present values. (B) is correct.
  },
}
